\IfFileExists{stacks-project.cls}{%
\documentclass{stacks-project}
}{%
\documentclass{amsart}
}

% For dealing with references we use the comment environment
\usepackage{verbatim}
\newenvironment{reference}{\comment}{\endcomment}
%\newenvironment{reference}{}{}
\newenvironment{slogan}{\comment}{\endcomment}
\newenvironment{history}{\comment}{\endcomment}

% For commutative diagrams we use Xy-pic
\usepackage[all]{xy}

% We use 2cell for 2-commutative diagrams.
\xyoption{2cell}
\UseAllTwocells

% We use multicol for the list of chapters between chapters
\usepackage{multicol}

% This is generally recommended for better output
\usepackage{lmodern}
\usepackage[T1]{fontenc}

% For cross-file-references
\usepackage{xr-hyper}

% Package for hypertext links:
\usepackage{hyperref}

% For any local file, say "hello.tex" you want to link to please
% use \externaldocument[hello-]{hello}
\externaldocument[introduction-]{introduction}
\externaldocument[conventions-]{conventions}
\externaldocument[sets-]{sets}
\externaldocument[categories-]{categories}
\externaldocument[topology-]{topology}
\externaldocument[sheaves-]{sheaves}
\externaldocument[sites-]{sites}
\externaldocument[stacks-]{stacks}
\externaldocument[fields-]{fields}
\externaldocument[algebra-]{algebra}
\externaldocument[brauer-]{brauer}
\externaldocument[homology-]{homology}
\externaldocument[derived-]{derived}
\externaldocument[simplicial-]{simplicial}
\externaldocument[more-algebra-]{more-algebra}
\externaldocument[smoothing-]{smoothing}
\externaldocument[modules-]{modules}
\externaldocument[sites-modules-]{sites-modules}
\externaldocument[injectives-]{injectives}
\externaldocument[cohomology-]{cohomology}
\externaldocument[sites-cohomology-]{sites-cohomology}
\externaldocument[dga-]{dga}
\externaldocument[dpa-]{dpa}
\externaldocument[sdga-]{sdga}
\externaldocument[hypercovering-]{hypercovering}
\externaldocument[schemes-]{schemes}
\externaldocument[constructions-]{constructions}
\externaldocument[properties-]{properties}
\externaldocument[morphisms-]{morphisms}
\externaldocument[coherent-]{coherent}
\externaldocument[divisors-]{divisors}
\externaldocument[limits-]{limits}
\externaldocument[varieties-]{varieties}
\externaldocument[topologies-]{topologies}
\externaldocument[descent-]{descent}
\externaldocument[perfect-]{perfect}
\externaldocument[more-morphisms-]{more-morphisms}
\externaldocument[flat-]{flat}
\externaldocument[groupoids-]{groupoids}
\externaldocument[more-groupoids-]{more-groupoids}
\externaldocument[etale-]{etale}
\externaldocument[chow-]{chow}
\externaldocument[intersection-]{intersection}
\externaldocument[pic-]{pic}
\externaldocument[weil-]{weil}
\externaldocument[adequate-]{adequate}
\externaldocument[dualizing-]{dualizing}
\externaldocument[duality-]{duality}
\externaldocument[discriminant-]{discriminant}
\externaldocument[derham-]{derham}
\externaldocument[local-cohomology-]{local-cohomology}
\externaldocument[algebraization-]{algebraization}
\externaldocument[curves-]{curves}
\externaldocument[resolve-]{resolve}
\externaldocument[models-]{models}
\externaldocument[functors-]{functors}
\externaldocument[equiv-]{equiv}
\externaldocument[pione-]{pione}
\externaldocument[etale-cohomology-]{etale-cohomology}
\externaldocument[proetale-]{proetale}
\externaldocument[relative-cycles-]{relative-cycles}
\externaldocument[more-etale-]{more-etale}
\externaldocument[trace-]{trace}
\externaldocument[crystalline-]{crystalline}
\externaldocument[spaces-]{spaces}
\externaldocument[spaces-properties-]{spaces-properties}
\externaldocument[spaces-morphisms-]{spaces-morphisms}
\externaldocument[decent-spaces-]{decent-spaces}
\externaldocument[spaces-cohomology-]{spaces-cohomology}
\externaldocument[spaces-limits-]{spaces-limits}
\externaldocument[spaces-divisors-]{spaces-divisors}
\externaldocument[spaces-over-fields-]{spaces-over-fields}
\externaldocument[spaces-topologies-]{spaces-topologies}
\externaldocument[spaces-descent-]{spaces-descent}
\externaldocument[spaces-perfect-]{spaces-perfect}
\externaldocument[spaces-more-morphisms-]{spaces-more-morphisms}
\externaldocument[spaces-flat-]{spaces-flat}
\externaldocument[spaces-groupoids-]{spaces-groupoids}
\externaldocument[spaces-more-groupoids-]{spaces-more-groupoids}
\externaldocument[bootstrap-]{bootstrap}
\externaldocument[spaces-pushouts-]{spaces-pushouts}
\externaldocument[spaces-chow-]{spaces-chow}
\externaldocument[groupoids-quotients-]{groupoids-quotients}
\externaldocument[spaces-more-cohomology-]{spaces-more-cohomology}
\externaldocument[spaces-simplicial-]{spaces-simplicial}
\externaldocument[spaces-duality-]{spaces-duality}
\externaldocument[formal-spaces-]{formal-spaces}
\externaldocument[restricted-]{restricted}
\externaldocument[spaces-resolve-]{spaces-resolve}
\externaldocument[formal-defos-]{formal-defos}
\externaldocument[defos-]{defos}
\externaldocument[cotangent-]{cotangent}
\externaldocument[examples-defos-]{examples-defos}
\externaldocument[algebraic-]{algebraic}
\externaldocument[examples-stacks-]{examples-stacks}
\externaldocument[stacks-sheaves-]{stacks-sheaves}
\externaldocument[criteria-]{criteria}
\externaldocument[artin-]{artin}
\externaldocument[quot-]{quot}
\externaldocument[stacks-properties-]{stacks-properties}
\externaldocument[stacks-morphisms-]{stacks-morphisms}
\externaldocument[stacks-limits-]{stacks-limits}
\externaldocument[stacks-cohomology-]{stacks-cohomology}
\externaldocument[stacks-perfect-]{stacks-perfect}
\externaldocument[stacks-introduction-]{stacks-introduction}
\externaldocument[stacks-more-morphisms-]{stacks-more-morphisms}
\externaldocument[stacks-geometry-]{stacks-geometry}
\externaldocument[moduli-]{moduli}
\externaldocument[moduli-curves-]{moduli-curves}
\externaldocument[examples-]{examples}
\externaldocument[exercises-]{exercises}
\externaldocument[guide-]{guide}
\externaldocument[desirables-]{desirables}
\externaldocument[coding-]{coding}
\externaldocument[obsolete-]{obsolete}
\externaldocument[fdl-]{fdl}
\externaldocument[index-]{index}

% Theorem environments.
%
\theoremstyle{plain}
\newtheorem{theorem}[subsection]{Theorem}
\newtheorem{proposition}[subsection]{Proposition}
\newtheorem{lemma}[subsection]{Lemma}

\theoremstyle{definition}
\newtheorem{definition}[subsection]{Definition}
\newtheorem{example}[subsection]{Example}
\newtheorem{exercise}[subsection]{Exercise}
\newtheorem{situation}[subsection]{Situation}

\theoremstyle{remark}
\newtheorem{remark}[subsection]{Remark}
\newtheorem{remarks}[subsection]{Remarks}

\numberwithin{equation}{subsection}

% Macros
%
\def\lim{\mathop{\mathrm{lim}}\nolimits}
\def\colim{\mathop{\mathrm{colim}}\nolimits}
\def\Spec{\mathop{\mathrm{Spec}}}
\def\Hom{\mathop{\mathrm{Hom}}\nolimits}
\def\Ext{\mathop{\mathrm{Ext}}\nolimits}
\def\SheafHom{\mathop{\mathcal{H}\!\mathit{om}}\nolimits}
\def\SheafExt{\mathop{\mathcal{E}\!\mathit{xt}}\nolimits}
\def\Sch{\mathit{Sch}}
\def\Mor{\mathop{\mathrm{Mor}}\nolimits}
\def\Ob{\mathop{\mathrm{Ob}}\nolimits}
\def\Sh{\mathop{\mathit{Sh}}\nolimits}
\def\NL{\mathop{N\!L}\nolimits}
\def\CH{\mathop{\mathrm{CH}}\nolimits}
\def\proetale{{pro\text{-}\acute{e}tale}}
\def\etale{{\acute{e}tale}}
\def\QCoh{\mathit{QCoh}}
\def\Ker{\mathop{\mathrm{Ker}}}
\def\Im{\mathop{\mathrm{Im}}}
\def\Coker{\mathop{\mathrm{Coker}}}
\def\Coim{\mathop{\mathrm{Coim}}}

% Boxtimes
%
\DeclareMathSymbol{\boxtimes}{\mathbin}{AMSa}{"02}

%
% Macros for moduli stacks/spaces
%
\def\QCohstack{\mathcal{QC}\!\mathit{oh}}
\def\Cohstack{\mathcal{C}\!\mathit{oh}}
\def\Spacesstack{\mathcal{S}\!\mathit{paces}}
\def\Quotfunctor{\mathrm{Quot}}
\def\Hilbfunctor{\mathrm{Hilb}}
\def\Curvesstack{\mathcal{C}\!\mathit{urves}}
\def\Polarizedstack{\mathcal{P}\!\mathit{olarized}}
\def\Complexesstack{\mathcal{C}\!\mathit{omplexes}}
% \Pic is the operator that assigns to X its picard group, usage \Pic(X)
% \Picardstack_{X/B} denotes the Picard stack of X over B
% \Picardfunctor_{X/B} denotes the Picard functor of X over B
\def\Pic{\mathop{\mathrm{Pic}}\nolimits}
\def\Picardstack{\mathcal{P}\!\mathit{ic}}
\def\Picardfunctor{\mathrm{Pic}}
\def\Deformationcategory{\mathcal{D}\!\mathit{ef}}
\usepackage{cleveref}
\crefname{originalreference}{original reference}{original references}
\Crefname{originalreference}{Original reference}{Original references}
\makeatletter
\newlabel{algebraic-lemma-characterize-representable-by-space}{{02ZX}{source}{}{sourcepointer-1}{}}
\newlabel{algebraic-lemma-characterize-representable-by-space@cref}{{[originalreference][0][]02ZX}{[1][0][]source}{}{}{}}
\makeatother
\crefname{originalreference}{original reference}{original references}
\Crefname{originalreference}{Original reference}{Original references}
\makeatletter
\newlabel{algebraic-section-representable-by-algebraic-spaces}{{02ZV}{source}{}{sourcepointer-2}{}}
\newlabel{algebraic-section-representable-by-algebraic-spaces@cref}{{[originalreference][0][]02ZV}{[1][0][]source}{}{}{}}
\makeatother
\crefname{originalreference}{original reference}{original references}
\Crefname{originalreference}{Original reference}{Original references}
\makeatletter
\newlabel{algebraic-section-split}{{04SU}{source}{}{sourcepointer-3}{}}
\newlabel{algebraic-section-split@cref}{{[originalreference][0][]04SU}{[1][0][]source}{}{}{}}
\makeatother
\crefname{originalequation}{original equation}{original equations}
\Crefname{originalequation}{Original equation}{Original equations}
\makeatletter
\newlabel{equation-limit-preserving}{{06CU}{source}{}{sourcepointer-4}{}}
\newlabel{equation-limit-preserving@cref}{{[originalequation][0][]06CU}{[1][0][]source}{}{}{}}
\makeatother
\crefname{originalreference}{original reference}{original references}
\Crefname{originalreference}{Original reference}{Original references}
\makeatletter
\newlabel{limits-lemma-descend-finite-locally-free}{{06AC}{source}{}{sourcepointer-5}{}}
\newlabel{limits-lemma-descend-finite-locally-free@cref}{{[originalreference][0][]06AC}{[1][0][]source}{}{}{}}
\makeatother
\crefname{originalreference}{original reference}{original references}
\Crefname{originalreference}{Original reference}{Original references}
\makeatletter
\newlabel{limits-lemma-descend-finite-presentation}{{01ZM}{source}{}{sourcepointer-6}{}}
\newlabel{limits-lemma-descend-finite-presentation@cref}{{[originalreference][0][]01ZM}{[1][0][]source}{}{}{}}
\makeatother
\crefname{originalreference}{original reference}{original references}
\Crefname{originalreference}{Original reference}{Original references}
\makeatletter
\newlabel{spaces-limits-proposition-characterize-locally-finite-presentation}{{04AK}{source}{}{sourcepointer-7}{}}
\newlabel{spaces-limits-proposition-characterize-locally-finite-presentation@cref}{{[originalreference][0][]04AK}{[1][0][]source}{}{}{}}
\makeatother
\title{The finite Hilbert stack of a representable morphism is limit preserving on objects}
\author{The Stacks Project Authors}
\date{stacks-project:criteria.tex:local-snapshot-20261012; GNU FDL 1.2 or later}
\begin{document}
\maketitle
\noindent\textit{Editorial note.} No omission markers in-span; proof complete as written.\par\medskip
\noindent\textit{Allowed standard prerequisites (editorial).}\begin{itemize}
\item Descent of finite locally free morphisms. limits-lemma-descend-finite-presentation, limits-lemma-descend-finite-locally-free (cross-chapter, named). Descends the degree-d cover Z -> U to some U\_\{i'\}.
\item Limit characterization of lfp. spaces-limits-proposition-characterize-locally-finite-presentation (cross-chapter, named). Descends the morphism f\_xi to f\_\{i'\} after enlarging the index.
\item Representable morphisms as functors. algebraic-lemma-characterize-representable-by-space (cross-chapter, named). Translates objects of the 2-fibre product into morphisms.
\end{itemize}\medskip
\ifdefined\hypertarget\hypertarget{sourcepointer-1}{}\fi
\noindent\textit{Source reference.} Original reference 02ZX: focused-seeds/geometry/S9-crit-limit-preserving/ref-tags.txt, source lines 5--5. Navigation pointer only: the counted proof cites this statement by \textbackslash{}ref as a standing input; the deduction the counted proof has to make is carried out inside the counted span, and no part of that deduction is delegated to this pointer.\par\medskip
\ifdefined\hypertarget\hypertarget{sourcepointer-2}{}\fi
\noindent\textit{Source reference.} Original reference 02ZV: focused-seeds/geometry/S9-crit-limit-preserving/ref-tags.txt, source lines 6--6. Navigation pointer only: the counted proof cites this statement by \textbackslash{}ref as a standing input; the deduction the counted proof has to make is carried out inside the counted span, and no part of that deduction is delegated to this pointer.\par\medskip
\ifdefined\hypertarget\hypertarget{sourcepointer-3}{}\fi
\noindent\textit{Source reference.} Original reference 04SU: focused-seeds/geometry/S9-crit-limit-preserving/ref-tags.txt, source lines 7--7. Navigation pointer only: the counted proof cites this statement by \textbackslash{}ref as a standing input; the deduction the counted proof has to make is carried out inside the counted span, and no part of that deduction is delegated to this pointer.\par\medskip
\ifdefined\hypertarget\hypertarget{sourcepointer-4}{}\fi
\noindent\textit{Source reference.} Original equation 06CU: sources/stacks-criteria/pinned-original-source.tex, source lines 570--570. Navigation pointer only: the counted proof cites this statement by \textbackslash{}ref as a standing input; the deduction the counted proof has to make is carried out inside the counted span, and no part of that deduction is delegated to this pointer.\par\medskip
\ifdefined\hypertarget\hypertarget{sourcepointer-5}{}\fi
\noindent\textit{Source reference.} Original reference 06AC: focused-seeds/geometry/S9-crit-limit-preserving/ref-tags.txt, source lines 8--8. Navigation pointer only: the counted proof cites this statement by \textbackslash{}ref as a standing input; the deduction the counted proof has to make is carried out inside the counted span, and no part of that deduction is delegated to this pointer.\par\medskip
\ifdefined\hypertarget\hypertarget{sourcepointer-6}{}\fi
\noindent\textit{Source reference.} Original reference 01ZM: focused-seeds/geometry/S9-crit-limit-preserving/ref-tags.txt, source lines 9--9. Navigation pointer only: the counted proof cites this statement by \textbackslash{}ref as a standing input; the deduction the counted proof has to make is carried out inside the counted span, and no part of that deduction is delegated to this pointer.\par\medskip
\ifdefined\hypertarget\hypertarget{sourcepointer-7}{}\fi
\noindent\textit{Source reference.} Original reference 04AK: focused-seeds/geometry/S9-crit-limit-preserving/ref-tags.txt, source lines 10--10. Navigation pointer only: the counted proof cites this statement by \textbackslash{}ref as a standing input; the deduction the counted proof has to make is carried out inside the counted span, and no part of that deduction is delegated to this pointer.\par\medskip
\begin{lemma}
\label{lemma-limit-preserving}
Let $F : \mathcal{X} \to \mathcal{Y}$ be a $1$-morphism of stacks in groupoids
over $(\Sch/S)_{fppf}$. Assume $F$ is representable by algebraic
spaces and locally of finite presentation. Then
$$
p : \mathcal{H}_d(\mathcal{X}/\mathcal{Y}) \to \mathcal{Y}
$$
is limit preserving on objects.
\end{lemma}

\begin{proof}
This means we have to show the following: Given
\begin{enumerate}
\item an affine scheme $U = \lim_i U_i$ which is written as the
directed limit of affine schemes $U_i$ over $S$,
\item an object $y_i$ of $\mathcal{Y}$ over $U_i$ for some $i$, and
\item an object $\Xi = (U, Z, y, x, \alpha)$ of
$\mathcal{H}_d(\mathcal{X}/\mathcal{Y})$
over $U$ such that $y = y_i|_U$,
\end{enumerate}
then there exists an $i' \geq i$ and an object
$\Xi_{i'} = (U_{i'}, Z_{i'}, y_{i'}, x_{i'}, \alpha_{i'})$ of
$\mathcal{H}_d(\mathcal{X}/\mathcal{Y})$ over $U_{i'}$ with
$\Xi_{i'}|_U = \Xi$ and $y_{i'} = y_i|_{U_{i'}}$.
Namely, the last two equalities will take care of the commutativity of
(\ref{equation-limit-preserving}).

\medskip\noindent
Let $X_{y_i} \to U_i$ be an algebraic space representing the $2$-fibre
product
$$
(\Sch/U_i)_{fppf} \times_{y_i, \mathcal{Y}, F} \mathcal{X}.
$$
Note that $X_{y_i} \to U_i$ is locally of finite presentation by our
assumption on $F$. Write $\Xi $. It is clear that
$\xi = (Z, Z \to U_i, x, \alpha)$ is an object of the $2$-fibre product
displayed above, hence $\xi$ gives rise to a morphism
$f_\xi : Z \to X_{y_i}$ of algebraic spaces over $U_i$
(since $X_{y_i}$ is the functor of isomorphisms classes of objects of
$(\Sch/U_i)_{fppf} \times_{y, \mathcal{Y}, F} \mathcal{X}$, see
Algebraic Stacks,
Lemma \ref{algebraic-lemma-characterize-representable-by-space}).
By
Limits, Lemmas \ref{limits-lemma-descend-finite-presentation} and
\ref{limits-lemma-descend-finite-locally-free}
there exists an $i' \geq i$ and a finite locally free morphism
$Z_{i'} \to U_{i'}$ of degree $d$ whose base change to $U$ is $Z$. By
Limits of Spaces, Proposition
\ref{spaces-limits-proposition-characterize-locally-finite-presentation}
we may, after replacing $i'$ by a bigger index, assume there exists a
morphism $f_{i'} : Z_{i'} \to X_{y_i}$ such that
$$
\xymatrix{
Z \ar[d] \ar[r] \ar@/^3ex/[rr]^{f_\xi} &
Z_{i'} \ar[d] \ar[r]_{f_{i'}} & X_{y_i} \ar[d] \\
U \ar[r] & U_{i'} \ar[r] & U_i
}
$$
is commutative. We set
$\Xi_{i'} = (U_{i'}, Z_{i'}, y_{i'}, x_{i'}, \alpha_{i'})$
where
\begin{enumerate}
\item $y_{i'}$ is the object of $\mathcal{Y}$ over $U_{i'}$
which is the pullback of $y_i$ to $U_{i'}$,
\item $x_{i'}$ is the object of $\mathcal{X}$ over $Z_{i'}$ corresponding
via the $2$-Yoneda lemma to the $1$-morphism
$$
(\Sch/Z_{i'})_{fppf} \to
\mathcal{S}_{X_{y_i}} \to
(\Sch/U_i)_{fppf} \times_{y_i, \mathcal{Y}, F} \mathcal{X} \to
\mathcal{X}
$$
where the middle arrow is the equivalence which defines $X_{y_i}$
(notation as in
Algebraic Stacks, Sections
\ref{algebraic-section-representable-by-algebraic-spaces} and
\ref{algebraic-section-split}).
\item $\alpha_{i'} : y_{i'}|_{Z_{i'}} \to F(x_{i'})$ is the isomorphism
coming from the $2$-commutativity of the diagram
$$
\xymatrix{
(\Sch/Z_{i'})_{fppf} \ar[r] \ar[rd] &
(\Sch/U_i)_{fppf} \times_{y_i, \mathcal{Y}, F} \mathcal{X}
\ar[r] \ar[d] &
\mathcal{X} \ar[d]^F \\
& (\Sch/U_{i'})_{fppf} \ar[r] & \mathcal{Y}
}
$$
\end{enumerate}
Recall that $f_\xi : Z \to X_{y_i}$ was the morphism corresponding to
the object $\xi = (Z, Z \to U_i, x, \alpha)$ of
$(\Sch/U_i)_{fppf} \times_{y_i, \mathcal{Y}, F} \mathcal{X}$
over $Z$. By construction $f_{i'}$ is the morphism corresponding to
the object $\xi_{i'} = (Z_{i'}, Z_{i'} \to U_i, x_{i'}, \alpha_{i'})$.
As $f_\xi = f_{i'} \circ (Z \to Z_{i'})$ we see that
the object $\xi_{i'} = (Z_{i'}, Z_{i'} \to U_i, x_{i'}, \alpha_{i'})$ pulls
back to $\xi$ over $Z$. Thus $x_{i'}$ pulls back to $x$ and $\alpha_{i'}$
pulls back to $\alpha$. This means that $\Xi_{i'}$ pulls back
to $\Xi$ over $U$ and we win.
\end{proof}

\end{document}
